\chapter{Arbitrering}
\label{ch:arbitrering}

{\large\itshape Hur en buss utan mästare avgör vem som får sända.\par}

\begin{chapterbox}{I det här kapitlet}
\begin{itemize}
  \item Multimaster: ingen nod fördelar bussen, och ingen behöver göra det.
  \item Arbitreringen bit för bit, arbetad för hand med två och tre noder.
  \item Varför den kallas icke-förstörande, och vad det sparar.
  \item Prioritetsinversion, svält, och hur ett verkligt system undviker dem.
\end{itemize}
\end{chapterbox}

\section{Ingen mästare}

De flesta bussar löser åtkomsten med en mästare. En nod bestämmer vem som får prata, och de andra
väntar på tur. Det fungerar, och det kostar en nod som alla är beroende av.

CAN har ingen sådan. Varje nod får börja sända så fort bussen varit ledig i tre bitar
(intermission). Börjar två noder samtidigt löser bussen själv upp det, medan det händer, utan
att någon förlorar tid.

Mekanismen är den asymmetri \kapref{ch:buss} lade fram, använd på ett enda ställe: under
arbitreringsfältet jämför varje sändande nod, bit för bit, det den sänder med det som faktiskt
ligger på bussen.

\begin{aside}
\textbf{Regeln, i sin helhet:} sänder en nod \rec{} och läser tillbaka \dom{}, har den förlorat.
Den slutar sända omedelbart och blir mottagare.
\end{aside}

Sänder den dominant och läser dominant är allt som det ska --- antingen är den ensam, eller så
sänder någon annan samma bit. Sänder den recessivt och läser recessivt likaså.

Den enda kombinationen som säger något är alltså recessiv ut, dominant in: någon annan sänder,
och den andra har en nolla där vi har en etta. Eftersom identifieraren sänds MSB först betyder
det att den andras identifierare är numeriskt \emph{lägre}, och därmed har högre prioritet.

\section{Två noder, bit för bit}

Nod A vill sända identifieraren \code{0x2A7}, nod B \code{0x2A3}. Båda börjar samtidigt.

Binärt, elva bitar, MSB först:

\begin{console}
bit:        1 2 3 4 5 6 7 8 9 10 11
A (0x2A7):  0 1 0 1 0 1 0 0 1  1  1
B (0x2A3):  0 1 0 1 0 1 0 0 0  1  1
                            ^ här skiljer de sig
\end{console}

\begin{booktable}{@{}lllll@{}}
\thead{Bit} & \thead{A sänder} & \thead{B sänder} & \thead{Bussen} & \thead{Händelse} \\ \midrule
1--8 & samma & samma & samma & Ingen märker något \\
9 & \rec{} (1) & \dom{} (0) & \dom{} & A läser tillbaka dominant: \textbf{A förlorar} \\
10--11 & --- & sänder vidare & B:s bitar & B fortsätter ostört \\
\end{booktable}

Notera vad som \emph{inte} händer: ingen kollision, ingen förlorad tid, ingen omstart. B:s frame
fortsätter från bit 9 som om A aldrig hade funnits. A:s frame sändes aldrig --- den bara slutade,
mitt i, och A väntar ut B:s frame för att försöka igen efteråt.

Därför heter det \textbf{icke-förstörande arbitrering}: konflikten löses utan att någons data
förstörs.

\bookfigure{arbitration}{Arbitrering mellan två noder. A drar sig ur och blir
mottagare vid bit 9; B:s frame löper vidare oavbrutet.}{fig:arbitration}

\section{Tre noder}

Med fler noder är det samma regel, tillämpad samtidigt av alla. Tre noder med identifierarna
\code{0x100}, \code{0x0FF} och \code{0x2A7}:

\begin{console}
bit:        1 2 3 4 5 6 7 8 9 10 11
A (0x100):  0 0 1 0 0 0 0 0 0  0  0
B (0x0FF):  0 0 0 1 1 1 1 1 1  1  1
C (0x2A7):  0 1 0 1 0 1 0 0 1  1  1
              ^ ^
              | bit 3: A sänder recessivt, B dominant -> A faller bort
              bit 2: C sänder recessivt, A och B dominant -> C faller bort
\end{console}

Vid bit 1 sänder alla tre dominant, så ingen märker något. Vid bit 2 faller C bort, vid bit 3
faller A bort. B vinner, och har lägst identifierare --- vilket alltid är utfallet, eftersom
"lägst numeriskt" och "först dominant bit" är samma sak när bitarna sänds MSB först.

\section{Vad förloraren gör}

Den förlorande noden blir mottagare för resten av framen. Det är inte bara artighet: den
\emph{ska} ta emot den vinnande framen, eftersom den mycket väl kan vara adresserad till den.

När bussen blivit ledig igen försöker den om. Och eftersom dess identifierare är oförändrad
kommer den att förlora igen om samma högprioriterade nod sänder igen --- vilket leder till nästa
avsnitt.

\begin{simplification}
\note[Förenkling] Den förenklade kontrollern i \kapref{ch:kursen} tar \emph{inte} emot framen den
förlorade mot. Den lämnar sändarrollen, sätter sin felflagga och går tillbaka till vila. Det är
en medveten förenkling, och den är synlig för mjukvaran: en nod som förlorat en arbitrering
missar den frame som vann.
\end{simplification}

\section{Svält och prioritetsinversion}

Arbitreringen är rättvis mot prioriteten, inte mot noderna. En nod med hög identifierare på en
hårt belastad buss kan i princip aldrig komma fram --- den förlorar varje gång.

Det är ett känt och accepterat beteende, och det hanteras i systemdesignen snarare än i
protokollet:

\begin{itemize}
  \item \textbf{Bussbelastningen hålls nere.} Tumregeln i fordonssystem är att inte planera för
    mer än omkring 30--50~\% belastning, just för att lämna utrymme åt lågprioriterade frames.
  \item \textbf{Identifierarna planeras.} Prioritetsordningen sätts av vad som är tidskritiskt,
    inte av vad som råkade få vilket nummer.
  \item \textbf{Sändningar periodiseras.} En nod som sänder så ofta den kan är en nod som svälter
    ut alla under sig.
\end{itemize}

\begin{aside}
\note En frame som förlorat arbitreringen försenas, men blir inte fel. Det är därför CAN kallas
\emph{deterministiskt} för den högst prioriterade framen: dess värstafallsfördröjning går att
räkna ut. För övriga frames går det också, men svaret blir större ju längre ned i ordningen de
ligger.
\end{aside}

\section{Vad sändaren inte får veta}

En nod som förlorat arbitreringen vet att den förlorade. En nod som vann vet bara att den kom
igenom.

Framför allt: \textbf{sändaren får aldrig veta om någon förlorade mot den.} Det finns ingen
räknare, ingen rapport, ingen skillnad i framen. Det spelar ingen roll för CAN, men det spelar
roll för den som felsöker ett system där en nod aldrig verkar komma fram --- symptomet syns bara
hos förloraren.

\section*{Övningar}

\exercise{Två noder}{Räkning}
\label{ex:4:1}
Nod A sänder identifieraren \code{0x155}, nod B sänder \code{0x154}. Båda börjar samtidigt.
\begin{enumerate}[a)]
  \item Skriv upp båda identifierarna binärt, elva bitar, MSB först.
  \item Vid vilken bit skiljer de sig?
  \item Vem vinner, och vad gör förloraren?
\end{enumerate}

\exercise{Tre noder}{Räkning}
\label{ex:4:2}
Tre noder börjar sända samtidigt, med identifierarna \code{0x400}, \code{0x200} och \code{0x201}.
\begin{enumerate}[a)]
  \item Vid vilken bit faller den första noden bort?
  \item Vid vilken bit avgörs det mellan de två återstående?
  \item Vem vinner?
\end{enumerate}

\exercise{Samma identifierare}{Resonemang}
\label{ex:4:3}
Två noder sänder av misstag samma identifierare samtidigt, men med olika data.
\begin{enumerate}[a)]
  \item Vad händer under arbitreringsfältet?
  \item Vad händer när datafälten börjar skilja sig åt?
  \item Varför är det här ett konstruktionsfel snarare än ett driftfel?
\end{enumerate}

\exercise{Svält}{Resonemang}
\label{ex:4:4}
En nod med identifieraren \code{0x7FF} kommer aldrig fram på en hårt belastad buss.
\begin{enumerate}[a)]
  \item Varför inte?
  \item Är det ett fel i protokollet? Motivera.
  \item Ge två åtgärder på systemnivå som skulle hjälpa.
\end{enumerate}
